% Generated by scripts/assemble.py; corrections belong in source/corrections.json.
\examhead{11}{北京科技大学 2024 年数学分析试卷}{16}
% p017-b14; source PDF page 17
\par\Needspace{5\baselineskip}
\blockanchor{p017-b14}
\majorquestion{一}解答题. 每题 10 分, 共 40 分.\par
% p017-b15; source PDF page 17
\blockanchor{p017-b15}
\question{1}求极限 $\lim_{x\to 0}\left(\frac{1}{x^2} -\frac{1}{x\tan x}\right)$\par
% p017-b16; source PDF page 17
\blockanchor{p017-b16}
\question{2}求极限 $\lim_{x\to 0}\frac{1 + \frac{1}{2}x^2 - \sqrt{1 + x^2}}{(\cos x - e^{x^2})\sin x^2}.$\par
% p017-b17; source PDF page 17
\blockanchor{p017-b17}
\question{3}设 $g(0) = 0, g'(0) = 1$, 分析 $f(x) = \begin{cases} g(x) \sin \frac{1}{x}, & x > 0; \\ g(x) \cos x, & x \leq 0. \end{cases}$ 在 $x = 0$ 处的连续性和可导性.\par
% p017-b18; source PDF page 17
\blockanchor{p017-b18}
\question{4}求 $I = \int xe^{ax}\cos bx\mathrm{d}x$ 及 $J = \int xe^{ax}\sin bx\mathrm{d}x.$\par
% p017-b19; source PDF page 17
\par\Needspace{5\baselineskip}
\blockanchor{p017-b19}
\majorquestion{二}(10 分) 设正数列 $\{a_{n}\}$ 满足 $a_{n} = \frac{a_{n+1}^{2}}{n} + a_{n+1}, n = 1,2,\cdots$. 计算极限 $\lim_{n\to \infty}a_n\ln n$.\par
% p017-b20; source PDF page 17
\par\Needspace{5\baselineskip}
\par\Needspace{6\baselineskip}
\blockanchor{p017-b20}
\majorquestion{三}(10 分) 设 $f(x)$ 在 $[a, b]$ 上有二阶导数, 且 $f^{(k)}(a) = f^{(k)}(b) = 0 (k = 0, 1)$, 证明: 存在 $\xi \in (a, b)$, 使得\par
\nopagebreak[4]
% p017-b21; source PDF page 17
\blockanchor{p017-b21}
\displayformula{f ^ {(2)} (\xi) = f (\xi).}
% p018-b00; source PDF page 18
\par\Needspace{5\baselineskip}
\blockanchor{p018-b00}
\majorquestion{四}（10分）利用致密性定理证明闭区间上的连续函数必然是有界的.\par
% p018-b01; source PDF page 18
\par\Needspace{5\baselineskip}
\blockanchor{p018-b01}
\majorquestion{五}（15分）按照 $p$ 的范围来说明级数 $\sum_{n=1}^{\infty}\left(\frac{1}{n^p}-\ln\left(1+\frac{1}{n^p}\right)\right)(p>0)$ 的收敛性.\par
% p018-b02; source PDF page 18
\par\Needspace{5\baselineskip}
\blockanchor{p018-b02}
\majorquestion{六}(15 分) 求 $x = r \cos \varphi, y = r \sin \varphi, z = r \cot \alpha$ 在点 $M_0(r_0, \varphi_0)$ 处的切平面方程 ($\alpha$ 为常数).\par
% p018-b03; source PDF page 18
\par\Needspace{5\baselineskip}
\blockanchor{p018-b03}
\majorquestion{七}(10 分) 讨论级数 $\sum_{n=0}^{\infty} \int_{0}^{1} (-1)^{n}(1 - x)x^{n} \, \mathrm{d}x$ 的收敛性并计算其和.\par
% p018-b04; source PDF page 18
\par\Needspace{5\baselineskip}
\par\Needspace{6\baselineskip}
\blockanchor{p018-b04}
\majorquestion{八}（15分）设函数 $f(x)$ 连续， $\Sigma$ 是球面 $x^{2} + y^{2} + z^{2} = 1$ ，且 $a,b,c$ 是常数，证明：\par
\nopagebreak[4]
% p018-b05; source PDF page 18
\blockanchor{p018-b05}
\displayformula{\iint_ {\Sigma} f (a x + b y + c z) \mathrm{d} S = 2 \pi \int_ {- 1} ^ {1} f (\sqrt {a ^ {2} + b ^ {2} + c ^ {2}} u) \mathrm{d} u.}
% p018-b06; source PDF page 18
\par\Needspace{5\baselineskip}
\par\Needspace{6\baselineskip}
\blockanchor{p018-b06}
\majorquestion{九}(15 分) 设 $a > 0, b > 0, m \neq 0$ 为常数, 计算积分\par
\nopagebreak[4]
% p018-b07; source PDF page 18
\blockanchor{p018-b07}
\displayformula{I (a, b) = \int_ {0} ^ {+ \infty} \frac {e ^ {- a x} - e ^ {- b x}}{x} \cos (m x) d x.}
% p018-b08; source PDF page 18
\par\Needspace{5\baselineskip}
\par\Needspace{6\baselineskip}
\blockanchor{p018-b08}
\majorquestion{十}(10 分) 设 $P(x, y), Q(x, y)$ 在 $\mathbb{R}^2$ 上有一阶连续偏导数, 若对以任意点 $(x_0, y_0) \in \mathbb{R}^2$ 为中心, 以任意 $r > 0$ 为半径的上半圆周 $L_r: y - y_0 = \sqrt{r^2 - (x - x_0)^2}$, 均有\par
\nopagebreak[4]
% p018-b09; source PDF page 18
\blockanchor{p018-b09}
\displayformula{I (r) = \int_ {L _ {r}} P (x, y) \mathrm{d} x + Q (x, y) \mathrm{d} y = 0.}
% p018-b10; source PDF page 18
\blockanchor{p018-b10}
证明： $P(x,y)\equiv 0,\frac{\partial Q(x,y)}{\partial x}\equiv 0.$\par
